\documentclass[10pt]{article}
\usepackage[margin=20mm]{geometry}
\usepackage{amsmath,amssymb,fontspec,unicode-math,xurl,hyperref}
\IfFontExistsTF{DejaVu Serif}{\setmainfont{DejaVu Serif}}{\setmainfont{Latin Modern Roman}}
\setmathfont{latinmodern-math.otf}
\setlength{\parindent}{0pt}\setlength{\parskip}{6pt}\sloppy
\begin{document}
Abstract: Independent finite-state calculations identify state-ID collisions and dropped transition conditions in the pinned optimizer. The paper also contains a listener-routing pseudocode error, inconsistent ablation arithmetic, and an ML example whose final narrative is not supported by its displayed execution. Some issues are corrected in the later runtime, and the numerical corrections preserve the main ranking. The report separates these observations from unverified historical benchmark effects. Scientific adjudication remains pending.\par
Author: AI-assisted technical review\par
Date:\par
Title: Technical Verification: MetaAgent: Automatically Constructing Multi-Agent Systems Based on Finite State Machines\par
\# Paper identity and reviewed version\par
This report concerns MetaAgent: Automatically Constructing Multi-Agent Systems Based on Finite State Machines, by Yaolun Zhang et al., ICML 2025, PMLR volume 267; identifier `I25-43677`, OpenReview `vOxaD3hhPt`. The reviewed object is the frozen published proceedings PDF, retrieved on 25 October 2026. Page references below are one-based PDF pages, which coincide with the internal page numbers. They are not the continuous PMLR volume pagination.\par
\# Review scope and interpretation\par
MetaAgent generates agents, a finite control graph, natural-language transition conditions and listener lists, then asks an LLM to merge states. The contribution is an empirical framework, not a proved theorem about formal automata. Main Sections 1–5, Algorithms 1–2, the merging prompt, software criteria, example software FSM and selected ML deployment trace were inspected. Core execution and optimization source was read statically at commit , dated 2 December 2025. That later public snapshot is not assumed identical to the historical experimental implementation. No author code, LLM calls or generated tool actions were executed.\par
The finite graph is a useful control abstraction. Agent memories, tool outputs and environment state are additional execution state; graph expressiveness alone does not prove generalization or correctness. In particular, role similarity and overlapping tools do not establish behavioral equivalence of states. The Section 3.2.3 assertion that merging preserves functionality should be treated as an empirical objective unless transition guards, memory effects and observable outputs are checked.\par
There is a positive invariant: if every accepted merge replaces two distinct states by one, Algorithm 1 terminates after at most \(n-1\) merges. Under finite-time LLM calls, at most \[\sum_{k=2}^{n}\binom{k}{2}=\binom{n+1}{3}\] pair tests suffice, including repeated scans. This is a termination and call-count bound, not a quality guarantee or a token-cost bound. The independent script verifies the combinatorial identity for several sizes.\par
\# Concrete failures of graph preservation in the pinned optimizer\par
Problem 1 (State merging can collide with an existing identifier and discard a guard).  The inspected optimizer does not preserve two basic graph invariants after an accepted merge: uniqueness of state IDs and preservation of distinct transition conditions.\par
In , lines 470–475 remove two states and then choose the new ID as the length of the remaining list. Consider IDs \(1,2,3,4,5\), with state 1 initial, state 5 final, and states 2, 3 and 4 ordinary. If an allowed pair 2 and 4, assigned to different agents, is approved for merging, the surviving IDs are \(1,3,5\) and the new ID is also 3. The flags check at lines 119–121 does not exclude this case. The later renumbering dictionary at lines 515–533 maps the duplicated old ID 3 to the last new value, 4; edges to the original surviving state 3 are consequently redirected to the merged state rather than its own new ID 2. Renumbering after a collision cannot recover the lost distinction.\par
Independently, lines 332–340 deduplicate transitions using only \((\text{from},\text{to})\), ignoring their conditions. If merging produces two edges \(m\to f\) with guards \(A\) and \(B\), retaining only the first implements \(A\) rather than \(A\lor B\). The four-row Boolean enumeration finds a lost transition when \(A\) is false and \(B\) true. This failure can occur even if the state ID is chosen correctly and even when the LLM’s intended merged instruction is otherwise reasonable.\par
These are Major implementation defects in the pinned optimizer, conditional on the corresponding merge being chosen. The independent calculation operates on a five-label list and a two-guard toy graph; it neither imports the author module nor claims that the actual LLM selects this pair on a benchmark. A repair is to allocate a fresh ID before editing, assert uniqueness and valid edge endpoints, and preserve each condition or combine same-endpoint guards by a well-defined logical disjunction. The independent repaired examples retain both the surviving state and all four guard truth-table outcomes. They are not a general validation of semantic state merging. Original benchmark impact and a full patched implementation remain unverified.\par
\# Deployment algorithm and implementation do not describe the same trace\par
Problem 2 (Algorithm 2 routes output using the destination state’s listeners).  Algorithm 2, page 16, assigns \(s\leftarrow s_{\mathrm{target}}\) at line 10, then inserts the old state’s output into \(s.\mathrm{Lis}\) at lines 12–13. Sections 3.2–3.3 specify the originating state’s listeners instead.\par
Take \(A\to B\to C\) with listeners \(L(A)=\{B\}\) and \(L(B)=\{C\}\). After \(A\) generates its result and transitions to \(B\), the displayed algorithm sends that result to \(C\) and leaves \(B\) without it. A deterministic listener-map calculation verifies the discrepancy. Saving the source state before assignment and broadcasting through its listener list repairs this example.\par
Algorithm 2 also contains no final-state termination test despite the definition in Section 3.1. Its null-transition branch calls the agent a second time but does not verify that second output before the next iteration begins. These are material omissions for a literal implementation, although they may be pseudocode mistakes. The pinned , lines 140–148 and 173–177, already checks submission/final state and broadcasts through the source state’s listener list. Thus the listener issue is Minor for the paper’s pseudocode, with an observed implementation-level correction; it is not evidence of that particular runtime failure in the released code.\par
The same source implements transition choice by parsing a marker in the task agent’s own output (lines 101–105, 131 and 149), whereas the paper describes a separately queried condition verifier. It also uses a local four-conversation cap and a separate count of graph transitions. These distinctions require version-specific documentation; the paper’s algorithm cannot alone establish the behavior of this later release. No historical version reconciliation was performed.\par
\# Numerical claims and execution evidence\par
Problem 3 (Several reported summary statistics cannot be reconstructed from their displayed inputs).  Table 5 uses inconsistent averages for its foundation-model ablation, and Table 6 overstates one optimization reduction.\par
The five displayed dataset scores in Table 5 give the following arithmetic means:\par
Designer / executor | Printed | Mean of five entries\par
GPT-3.5 / GPT-3.5 | .12 | .146\par
GPT-3.5 / GPT-4o | .39 | .462\par
GPT-4o / GPT-3.5 | .26 | .306\par
GPT-4o / GPT-4o | .83 | .832\par
The first three printed values are consistent with dividing by six instead of five, up to rounding. The row labels and accompanying text also disagree: the paragraph attributes .26 and .35 to designer/executor replacement, respectively, while those are not the table’s corresponding two values. Correcting the means still supports the qualitative ordering that the weaker executor loses more than the weaker designer, but changes the magnitudes.\par
For software optimization, \((.85-.65)/.85=23.53\%\), not the stated \(35.3\%\) decrease. Likewise software checkpoint success rising from .35 to .85 is a 50 percentage-point increase, or \(142.86\%\) relative improvement, not literally “50\% more” checkpoints. These are Minor numerical/reporting errors: the main ordering survives. Table 3’s claim of roughly 97\% of Data Interpreter performance is consistent with \(.83/.86=96.51\%\). Table 7’s displayed design-plus-deployment totals also add correctly. Its software cost coverage is six tasks, however, while Table 4 and Appendix B name five; the sixth case and its relationship to the performance aggregate are unresolved.\par
Problem 4 (The provided ML example does not support its final model-selection narrative).  Appendix H.1 displays a random-forest run with accuracy .8324, followed by a report naming gradient boosting as best and listing several additional unshown models and features.\par
Pages 24–26 show preprocessing that drops the passenger-name column and trains a random-forest classifier. The recorded execution result on page 26 is .8324. The reporting agent then claims title extraction, five evaluated model families, random forest at .82 and gradient boosting best at .83. The displayed trace contains no execution evidence for these additions. This is a Major gap in using the example as evidence of reliable execution-grounded reporting; it is not evidence that all table scores were fabricated or that the random-forest .8324 result itself is wrong. The correct action is to bind reported metrics to saved predictions, evaluation outputs and model identities, and to disclose whether Table 3 came from those outputs or unconstrained generated summaries. The small wrapper only returns the system result and supplies no independent ML scorer. Other benchmark-specific evaluators or logs may exist outside this snapshot; their absence here is retained as an uncertainty.\par
\# Strengths, remaining tests and conclusion\par
The framework integrates practical tools and backtracking into a reusable task-level graph, publishes code, and supplies ablations and explicit software checkpoints. These are useful contributions. The reported advantages are package-level empirical comparisons: Table 6 gives tool-free GPQA performance .52 versus full-system .60, while Table 2 lists .55 for Self-Refine. Without equal tools and budgets, the headline gain cannot be attributed solely to the FSM or its verifier. The table entries alone do not establish paired statistical significance, and a zero sampling temperature does not establish identical service behavior across runs. Five software cases with four criteria each give coarse scores; one criterion, aesthetically displaying weather data, still requires judgment.\par
The supported findings concern graph-integrity failures in a specific later source commit, correctable pseudocode routing, reproducible arithmetic discrepancies, and a concrete mismatch between an execution trace and its narrative. Finite calculations validate the local counterexamples and repairs. They do not replace LLM deployment tests, evaluation-set isolation checks, comprehensive tool audits, independent scoring or a search for later corrections. The main empirical ordering is not overturned by the arithmetic corrections alone. Focused independent AI acceptance is complete. Adjudication, original-run impact and author adoption remain pending or unknown; there has been no human review.\par
1 Yaolun Zhang et al., MetaAgent: Automatically Constructing Multi-Agent Systems Based on Finite State Machines, PMLR 267 (2025), published proceedings version. https://proceedings.mlr.press/v267/zhang25bc.html.\par
\# Evidence and reproducibility\par
The exact source PDF has SHA-256\par
The matching `verification/0308` directory contains the source text extraction, independently written calculation script and actual output, inspected page images, coverage ledger, findings and provenance. Source code supplied by the authors was not executed. The calculation is an independent finite or analytic diagnostic, not a reproduction of the original training experiments. The final PDF is compiled from this TeX; the compiler log and visual receipt identify the exact bytes inspected. This initial review was performed by the delegated AI unit with GPT-6 Astra / Max formally selected. A separate AI unit subsequently checked the stated Problems against the frozen original source, assumptions, decisive arguments, counterevidence and impact boundaries. Its fresh evidence and exact artifact bindings are retained in the same verification directory. No human review has occurred; scientific adjudication is pending and author adoption is unknown.\par
\end{document}
